

%% 
%% Copyright 2019-2020 Elsevier Ltd
%% 
%% This file is part of the 'CAS Bundle'.
%% --------------------------------------
%% 
%% It may be distributed under the conditions of the LaTeX Project Public
%% License, either version 1.2 of this license or (at your option) any
%% later version.  The latest version of this license is in
%%    http://www.latex-project.org/lppl.txt
%% and version 1.2 or later is part of all distributions of LaTeX
%% version 1999/12/01 or later.
%% 
%% The list of all files belonging to the 'CAS Bundle' is
%% given in the file `manifest.txt'.
%% 
%% Template article for cas-sc documentclass for 
%% double column output.

%\documentclass[a4paper,fleqn,longmktitle]{cas-sc}
\documentclass[a4paper,fleqn]{cas-sc}

% \usepackage[numbers]{natbib}
%\usepackage[authoryear]{natbib}
\usepackage[authoryear,longnamesfirst]{natbib}
\usepackage{amssymb}
\usepackage{pifont}
\newcommand{\cmark}{\ding{51}}  % ✓
\newcommand{\xmark}{\ding{55}}  % ✗                
\usepackage{svg}
\usepackage{booktabs}         
\usepackage{array} 
\usepackage[table]{xcolor}   % colori nelle celle     
%% The amsmath package provides various useful equation environments.
\usepackage{amsmath}
%% The amsthm package provides extended theorem environments
%% \usepackage{amsthm}
%% The lineno packages adds line numbers. Start line numbering with
%% \begin{linenumbers}, end it with \end{linenumbers}. Or switch it on
%% for the whole article with \linenumbers.
\usepackage{lineno}
\usepackage{pdfpages}

\usepackage[utf8]{inputenc}
\usepackage[T1]{fontenc}

%\usepackage{array}
\usepackage{enumitem}
\usepackage{makecell}
%\usepackage{geometry}
%\usepackage{booktabs}
\usepackage{graphicx}
\usepackage{tabularx}
\usepackage{comment}
\usepackage{tabularray}
%\usepackage{pdflscape} 
\usepackage{multirow}
\usepackage{hhline}
\usepackage{rotating}
\usepackage{colortbl}
\usepackage{float}
\usepackage{subcaption}
\usepackage{pdfpages}
\usepackage{changepage}
%\usepackage{enumerate}
%\usepackage[table,xcdraw]{xcolor}  
\usepackage{url}
\usepackage[scaled]{helvet}
%\usepackage{color}
%\usepackage[colorlinks=true, linkcolor=black, urlcolor=black, citecolor=black, filecolor=black, backref]{hyperref}  

\usepackage{cancel}
% Beamer presentation requires \usepackage{colortbl} instead of
\usepackage[normalem]{ulem}
\usepackage{booktabs}
\usepackage{longtable}
\usepackage{multirow}
\usepackage{graphicx}  
\usepackage{booktabs,tabularx}


\newcommand{\DD}[2]{\textcolor{blue}{#1}[\textcolor{red}{#2}]}
\newcommand{\GM}[2]{\textcolor{red}{#1}[\textcolor{teal}{#2}]}
\newcommand{\NC}[2]{\textcolor{olive}{#1}[\textcolor{olive}{#2}]}
\newcommand{\SA}[2]{\textcolor{olive}{#1}[\textcolor{magenta}{#2}]}

%\newcommand{\PSRQone}{PS-RQ1: What kind and how many primary studies have been identified per publication year?}
\newcommand{\PSRQone}{PS-RQ1: How are the primary studies distributed by venue type and publication year?}
%\newcommand{\PSRQtwo}{PS-RQ2: What are the venues and areas of interest of the primary studies that have been published?}
\newcommand{\PSRQtwo}{PS-RQ2: What are the venues where the primary studies have been published?}
%\newcommand{\PSRQthree}{PS-RQ3: What are the authors’ affiliation countries of the selected primary studies?}
\newcommand{\PSRQthree}{PS-RQ3: How are primary studies distributed across authors’ affiliation countries?}

\newcommand{\RSRQone}{RS-RQ1: What are the adopted definitions of customer churn?}
\newcommand{\RSRQtwo}{RS-RQ2: What features have been used for customer churn prediction, and which have been found to be more effective?}
\newcommand{\RSRQthree}{RS-RQ3: How have the available features been pre-processed and/or engineered before analyzing them using AI techniques?}
\newcommand{\RSRQfour}{RS-RQ4: What AI techniques have been applied and found to be most effective in predicting customer churn in the retail sector?}
\newcommand{\RSRQfive}{RS-RQ5: What are the most commonly used metrics to evaluate the performance of churn prediction techniques?}
\newcommand{\RSRQsix}{RS-RQ6: Which publicly available datasets have been used for retail customer churn prediction, and what are their main characteristics?}
\newcommand{\RSRQseven}{RS-RQ7: What are the challenges and limitations reported in using AI techniques for churn prediction?}
\newcommand{\RSRQeight}{RS-RQ8: What are the emerging trends and future directions for research in this area?}

\definecolor{lightgray}{HTML}{EFEFEF}

\widowpenalty=10000
\clubpenalty=10000

%%%Author definitions
\def\tsc#1{\csdef{#1}{\textsc{\lowercase{#1}}\xspace}}
\tsc{WGM}
\tsc{QE}
\tsc{EP}
\tsc{PMS}
\tsc{BEC}
\tsc{DE}
%%%

% Uncomment and use as if needed
%\newtheorem{theorem}{Theorem}
%\newtheorem{lemma}[theorem]{Lemma}
%\newdefinition{rmk}{Remark}
%\newproof{pf}{Proof}
%\newproof{pot}{Proof of Theorem \ref{thm}}

\begin{document}
%\let\WriteBookmarks\relax
%\def\floatpagepagefraction{1}
%\def\textpagefraction{.001}

% Short title
\shorttitle{}

% Short author
\shortauthors{}

% Main title of the paper
\title [mode = title]{Artificial Intelligence for Customer Churn Prediction in Non-Contractual Retail Environments: A Systematic Literature Review}                      
% Title footnote mark
% eg: \tnotemark[1]
%\tnotemark[1,2]

% Title footnote 1.
% eg: \tnotetext[1]{Title footnote text}
% \tnotetext[<tnote number>]{<tnote text>} 
%\tnotetext[1]{This document is the results of the research
%   project funded by the National Science Foundation.}

%\tnotetext[2]{The second title footnote which is a longer text matter
%   to fill through the whole text width and overflow into
%   another line in the footnotes area of the first page.}


% First author
%
% Options: Use if required
% eg: \author[1,3]{Author Name}[type=editor,
%       style=chinese,
%       auid=000,
%       bioid=1,
%       prefix=Sir,
%       orcid=0000-0000-0000-0000,
%       facebook=<facebook id>,
%       twitter=<twitter id>,
%       linkedin=<linkedin id>,
%       gplus=<gplus id>]
\author[1,2]{Niccolò Ciucci Colli}[orcid=0009-0007-9049-6087]
% Corresponding author indication
\cormark[1]
% Footnote of the first author
%\fnmark[1]
% Email id of the first author
\ead{niccolo.ciuccicolli@uniroma1.it}
%  Credit authorship
\credit{Data curation, Formal analysis, Investigation, Methodology,Project administration, Visualization, Writing – original draft, Writing – review and editing}

% Second author
\author[1,2]{Syed Saad Saif}
%\fnmark[2]
\credit{Data curation, Formal analysis, Investigation,
        Visualization, Writing – original draft,
        Writing – review and editing}


% Third author
\author[2]{Giulio Maggiore}
%\fnmark[3]
\credit{Conceptualization, Investigation, Supervision,
        Validation, Writing – review and editing}

% Fourth author
\author[2]{Damiano Distante}
%\fnmark[4]
\credit{Conceptualization, Investigation, Methodology,
        Project administration, Supervision, Validation,
        Visualization, Writing – review and editing}

\affiliation[1]{organization={Department of Computer, Control, and Management Engineering, Sapienza University of Rome},
                addressline={Piazzale Aldo Moro 5},
                city={Rome},
                postcode={00185},
                country={Italy}}

\affiliation[2]{organization={Department of Law and Economics, University of Rome UnitelmaSapienza},
                addressline={Piazza Sassari 4},
                city={Rome},
                postcode={00161},
                country={Italy}}

% Corresponding author text
\cortext[cor1]{Corresponding author}

% Footnote text
%\fntext[fn1]{This is the first author footnote. but is common to third
%  author as well.}
%\fntext[fn2]{Another author footnote, this is a very long footnote and
%  it should be a really long footnote. But this footnote is not yet
%  sufficiently long enough to make two lines of footnote text.}



\maketitle

%\printcredits


\end{document}

